% Zumdahl Chapter 13 free-response set -- 2 long (10) + 3 short (4) = 32 pts.
\documentclass{apchem-exam}

\apcunitword{Chapter}
\apcunit{13}
\apcunittitle{Chemical Equilibrium}
\apcdoctitle{Free-Response Set}
\apcsubtitle{Zumdahl \S13.1--\S13.7 \textbullet\ 5 questions \textbullet\ 32 points \textbullet\ 50 minutes}

\begin{document}
\apctitleblock

\begin{apcnote}[Scope --- one CED exclusion applies]
Covers \textbf{CED 7.1--7.2} (\S13.1), \textbf{7.3} and \textbf{7.5}
(\S13.2, \S13.5), \textbf{7.4} (\S13.2--13.4), \textbf{7.6}
(\S13.3--13.4), \textbf{7.7} (\S13.6), \textbf{7.8} (\S13.1) and
\textbf{7.9--7.10} (\S13.7). Equilibrium sits in CED Unit 7 and is one of
the most heavily tested ideas on the exam.

The CED exclusion on topic 7.3 removes \textbf{converting between $K_c$
and $K_p$}. Zumdahl derives $K_p = K_c(RT)^{\Delta n}$ in \S13.3; that
conversion is not assessed and is not asked for here. Writing and using
$K_p$ directly from partial pressures remains examinable.
\end{apcnote}

\examsection{II}{Free Response}{5 questions \textbullet\ 32 points \textbullet\ 50 minutes}{\calculatorok}

% =====================================================================
\frq{long}{10}{Zumdahl \S13.6 \textbullet\ CED 7.4, 7.7 \textbullet\ an ICE table}

At a certain temperature, $K_c = 50.0$ for
\[ \ce{H2(g) + I2(g) <=> 2HI(g)} \]
A vessel is charged with \SI{1.00}{\Molar} \ce{H2} and
\SI{1.00}{\Molar} \ce{I2} and no \ce{HI}.

\begin{parts}
  \item Write the equilibrium constant expression.
        \answerlines[2]{$K_c = \dfrac{[\ce{HI}]^2}
        {[\ce{H2}][\ce{I2}]}$}
  \item Construct an ICE table for the reaction. \workspace[3.4cm]
        \keyonly{\begin{apcanswer}\begin{center}
        \begin{tabular}{lccc}
        \hline
         & \ce{H2} & \ce{I2} & \ce{HI} \\
        \hline
        Initial     & 1.00   & 1.00   & 0 \\
        Change      & $-x$   & $-x$   & $+2x$ \\
        Equilibrium & $1.00-x$ & $1.00-x$ & $2x$ \\
        \hline
        \end{tabular}
        \end{center}
        The \ce{HI} change is $+2x$, not $+x$ --- the coefficient
        matters.\end{apcanswer}}
  \item Calculate the equilibrium concentration of \ce{HI}.
        \workspace[3.4cm]
        \keyonly{\begin{apcanswer}$\dfrac{(2x)^2}{(1.00-x)^2} = 50.0$

        Both sides are perfect squares, so take the square root of each
        --- no quadratic formula is needed:

        $\dfrac{2x}{1.00-x} = \sqrt{50.0} = 7.071$

        $2x = 7.071 - 7.071x \;\Rightarrow\; 9.071x = 7.071
        \;\Rightarrow\; x = \num{0.7795}$

        $[\ce{HI}] = 2x = \mathbf{1.56}$~M, and
        $[\ce{H2}] = [\ce{I2}] = \num{0.220}$~M.

        Check: $(1.559)^2/(0.2205)^2 = 50.0$
        \checkmark\end{apcanswer}}
  \item Explain why the approximation $1.00 - x \approx 1.00$ would be
        invalid here. \answerlines[3]{That approximation is only safe when
        $x$ is small compared with the initial concentration, which
        happens when $K$ is very small. Here $K = 50.0$ is large, the
        reaction proceeds far toward products, and $x = 0.78$ is most of
        the initial \SI{1.00}{\Molar}. Dropping it would be a
        \SI{78}{\percent} error.}
\end{parts}

\begin{rubric}
  \rpoint{2}{(a) correct expression with \ce{HI} squared; 1 pt if the
    square is missing}
  \rpoint{2}{(b) 1 pt initial row; 1 pt the change row with $+2x$ for
    \ce{HI}. Using $+x$ caps the whole question at 5}
  \rpoint{3}{(c) 1 pt setting up $K$ with the ICE values; 1 pt solving
    (square-root method or quadratic); 1 pt $[\ce{HI}] = \SI{1.56}{\Molar}$}
  \rpoint{3}{(d) 1 pt the approximation needs $x$ small relative to
    initial; 1 pt $K$ is large here; 1 pt $x$ is a large fraction of
    \num{1.00}}
\end{rubric}

% =====================================================================
\frq{long}{10}{Zumdahl \S13.7 \textbullet\ CED 7.9--7.10 \textbullet\ Le Ch\^{a}telier and the reaction quotient}

For the exothermic equilibrium
\[ \ce{N2(g) + 3H2(g) <=> 2NH3(g)} \qquad \Delta H < 0 \]
state and justify the effect of each change on the amount of \ce{NH3} at
equilibrium.

\begin{parts}
  \item The volume of the container is decreased.
        \answerlines[3]{The amount of \ce{NH3} \textbf{increases}.
        Decreasing the volume raises all partial pressures; the system
        responds by shifting toward the side with fewer moles of gas.
        There are 4 moles of gas on the left and 2 on the right, so the
        equilibrium shifts right.}
  \item The temperature is increased.
        \answerlines[3]{The amount of \ce{NH3} \textbf{decreases}. The
        forward reaction is exothermic, so heat behaves as a product.
        Adding heat shifts the system in the endothermic direction --- to
        the left. Note this also \emph{changes the value of $K$}, which
        temperature alone among these changes does.}
  \item A catalyst is added.
        \answerlines[2]{\textbf{No change} in the amount of \ce{NH3}. A
        catalyst lowers the activation energy of the forward and reverse
        reactions equally, so equilibrium is reached sooner but its
        position, and $K$, are unaltered.}
  \item At one moment the mixture contains
        $[\ce{N2}] = \SI{0.10}{\Molar}$,
        $[\ce{H2}] = \SI{0.10}{\Molar}$ and
        $[\ce{NH3}] = \SI{0.50}{\Molar}$, and $K_c = 0.50$. Determine the
        direction of net change. \workspace[3cm]
        \keyonly{\begin{apcanswer}$Q = \dfrac{[\ce{NH3}]^2}
        {[\ce{N2}][\ce{H2}]^3} = \dfrac{(0.50)^2}{(0.10)(0.10)^3}
        = \dfrac{0.25}{1.0\times10^{-4}} = \num{2.5e3}$

        $Q \gg K$ (2500 against 0.50), so there is far too much product.
        The system shifts \textbf{to the left}, consuming \ce{NH3} until
        $Q$ falls to $K$.\end{apcanswer}}
\end{parts}

\begin{rubric}
  \rpoint{3}{(a) 1 pt increases; 1 pt shift toward fewer moles of gas;
    1 pt the 4-versus-2 count}
  \rpoint{3}{(b) 1 pt decreases; 1 pt exothermic so heat acts as a
    product; 1 pt noting $K$ itself changes with temperature}
  \rpoint{1}{(c) no change, with the reason}
  \rpoint{3}{(d) 1 pt correct $Q$ expression with the cube on \ce{H2};
    1 pt $Q = \num{2.5e3}$; 1 pt shifts left because $Q > K$. Omitting the
    cube caps this part at 1}
\end{rubric}

% =====================================================================
\frq{short}{4}{Zumdahl \S13.4 \textbullet\ CED 7.6 \textbullet\ heterogeneous equilibria}

\begin{parts}
  \item Write the equilibrium constant expression for
        \ce{CaCO3(s) <=> CaO(s) + CO2(g)}.
        \answerlines[2]{$K = [\ce{CO2}]$ (or $K_p = P_{\ce{CO2}}$).}
  \item Explain why the two solids are omitted.
        \answerlines[3]{A pure solid has a fixed density and composition,
        so its ``concentration'' does not change no matter how much is
        present. Adding or removing solid does not shift the equilibrium,
        so the solid is absorbed into the constant rather than written in
        the expression.}
\end{parts}

\begin{rubric}
  \rpoint{2}{(a) $K = [\ce{CO2}]$ alone; including either solid earns 0}
  \rpoint{2}{(b) 1 pt a pure solid has constant concentration; 1 pt its
    amount does not affect the position of equilibrium}
\end{rubric}

% =====================================================================
\frq{short}{4}{Zumdahl \S13.5 \textbullet\ CED 7.5 \textbullet\ the magnitude of K}

Two reactions at the same temperature have $K = 1\times10^{12}$ and
$K = 1\times10^{-8}$.

\begin{parts}
  \item State what each value indicates about the position of equilibrium.
        \answerlines[2]{$K = 10^{12}$: equilibrium lies far to the right;
        the mixture is essentially all products. $K = 10^{-8}$: far to the
        left; almost entirely reactants.}
  \item State whether a large $K$ means the reaction is fast, and explain.
        \answerlines[3]{No. $K$ describes only the \emph{position} of
        equilibrium --- how far the reaction goes --- and says nothing
        about how quickly it gets there. A reaction with a huge $K$ may be
        immeasurably slow if its activation energy is high. Extent and
        rate are independent.}
\end{parts}

\begin{rubric}
  \rpoint{2}{(a) 1 pt each interpretation}
  \rpoint{2}{(b) 1 pt ``no''; 1 pt $K$ gives extent, not rate}
\end{rubric}

% =====================================================================
\frq{short}{4}{Zumdahl \S13.1--\S13.2 \textbullet\ CED 7.1--7.2, 7.8 \textbullet\ the nature of equilibrium}

\begin{parts}
  \item Explain what is meant by describing chemical equilibrium as
        \emph{dynamic}. \answerlines[3]{Both the forward and reverse
        reactions continue at the molecular level, at equal rates. Nothing
        stops --- individual molecules keep reacting in both directions
        --- but because the rates are equal the concentrations no longer
        change.}
  \item State what a graph of concentration against time looks like once
        equilibrium is reached, and what it would look like if the
        reaction had simply stopped.
        \answerlines[3]{At equilibrium all curves become horizontal at
        non-zero values, with reactants still present. If the reaction had
        stopped instead, the curves would also flatten --- so the graph
        alone cannot distinguish the two. The evidence for dynamic
        behaviour comes from other experiments, such as isotopic labelling
        showing continued exchange.}
\end{parts}

\begin{rubric}
  \rpoint{2}{(a) 1 pt both reactions continue; 1 pt at equal rates so
    concentrations are constant}
  \rpoint{2}{(b) 1 pt horizontal lines at non-zero concentrations with
    reactants remaining; 1 pt recognizing the graph alone does not prove
    the process is dynamic}
\end{rubric}

\examtotal

\end{document}
